\documentclass[a4paper]{report}
\usepackage[magyar]{babel}
\usepackage[utf8]{inputenc}
\usepackage[T1]{fontenc}

\begin{document}

\textbf{Témavezető neve:} Szabó László Ferenc

\textbf{munkahelyének neve, tanszéke:} ELTE IK Algoritmusok és Alkalmazásaik Tanszék

\textbf{munkahelyének címe:} 1117 Budapest, Pázmány Péter sétány 1/C

\textbf{beosztás és iskolai végzettsége:} habilitált egyetemi docens

\

\textbf{A diplomamunka címe:} Cache-optimális algoritmusok elemzése

\textbf{A diplomamunka témája: (A témavezetővel konzultálva adja meg 1/2 - 1 oldal terjedelemben diplomamunka témájának  leírását )}

\

A diplomamunkám célja a cache-optimális és cache-oblivious algoritmusok bemutatása.
A nagy számításigényű problémák hatékony megoldása a mai számítógépeken a cache megfelelő használatán alapszik.
A cache-optimális algoritmusok a cache paramétereinek ismeretét feltételezik, míg a cache-oblivious algoritmusok univerzális megoldást próbálnak adnak.

A dolgozatban ismertetem a cache számítási modelljét, az ismert algoritmusokat és ezek elméleti elemzését.

Az elmélet kevés útmutatást ad ezen algoritmusok gyakorlati használatára.
A cache-optimális algoritmusok finomhangolást igényelnek, a cache-oblivious algoritmusok pedig arról hírhedtek, hogy a futási idejük többszöröse egy kézzel hangolt megoldásénak.
Az optimális megoldást nehezíti az elvárás, hogy egy egyszer megírt programot különböző, heterogén számítógépekből álló környezetben kívánják futtatni.

Az algoritmusok összehasonlíthatóságához szimulátort készítek.
A probléma összetettsége miatt arra koncentrálok, hogy meg lehessen adni
\begin{itemize}
\item érdekes bonyolultságú algoritmusokat,
\item érdekes méretű feladatokat,
\item ténylegesen előforduló többszintű cache-hierarchiákat,
\item és interferáló programokat.
\end{itemize}

\end{document}
